\section{Problem}
\label{sec:problem}

\paragraph{Cells.} A target is a boolean mask on a $W\times H$ lattice. Lattice position
$(x,y)$ is column $x$ and row $y$, and usable positions are the cells
$c_1,\dots,c_M$. Cell centers default to the integer coordinates, so one unit is one cell
pitch. Anisotropic pitches (e.g.\ for 4:3 thumbnails) scale the centers but leave the
lattice unchanged. A full rectangle, a rectangle whose last row is only partly usable,
a disk, a ring, and the shape of a letter are all masks.

\paragraph{Points and normalization.} Input points $q_1,\dots,q_N$ ($N\le M$) are mapped
into the cell frame before anything else. Unless stated otherwise, each axis of the
point bounding box is mapped affinely onto the bounding box of the usable cell centers
(\emph{bbox} normalization); $p_i$ denotes the mapped points. The mapping is part of the
problem: in every comparison, all methods receive identical $p_i$ and identical cell
centers.

\paragraph{Objective.} We seek an injective assignment $\sigma:\{1..N\}\to\{1..M\}$
minimizing
\begin{equation}
  C(\sigma) = \sum_{i=1}^N \lVert p_i - c_{\sigma(i)}\rVert^2 .
  \label{eq:cost}
\end{equation}
When $N=M$ this is a bijection. When $N<M$ the solver must also decide which $M-N$ cells
stay empty. Unlike the bijective case, the cell term $\sum_i \lVert c_{\sigma(i)}\rVert^2$
then depends on the choice, so the surplus case is not just a relabelling of the
bijective one.

\paragraph{A certificate of optimality.} Equation~\eqref{eq:cost} is the linear program
$\min \sum_{ij} c_{ij}x_{ij}$ with $\sum_j x_{ij}=1$ for every point,
$\sum_i x_{ij}\le1$ for every cell, and $x\ge0$, where $c_{ij}=\lVert p_i-c_j\rVert^2$.
Its constraint matrix is totally unimodular. We use the following elementary bound.

\begin{proposition}\label{prop:lb}
For any cell potentials $v\in\R^M$,
\[
  \mathrm{LB}(v) = \sum_{i=1}^N \min_{j} \bigl(c_{ij}-v_j\bigr) + \sum_{j=1}^M v_j - (M-N)\max_j v_j
\]
is a lower bound on the optimal cost. The bound is tight (equals the optimum) for
optimal dual potentials.
\end{proposition}
\begin{proof}
Let $\bar v=\max_j v_j$ and $v'_j=v_j-\bar v\le0$. Let $u_i=\min_j(c_{ij}-v'_j)$. Then
$(u,v')$ is feasible for the dual $\max\sum_i u_i+\sum_j v'_j$ subject to
$u_i+v'_j\le c_{ij}$ and $v'\le0$, and its objective equals $\mathrm{LB}(v)$ (shifting all
potentials by $\bar v$ changes every $u_i$ by $+\bar v$ and the cell sum by $-M\bar v$).
Weak duality gives the bound. Strong duality of the LP gives tightness.
\end{proof}

The bound needs a minimum over \emph{all} $M$ cells for every point. It does not depend
on how $v$ was obtained, so $C(\sigma)-\mathrm{LB}(v)$ certifies the gap of any
assignment $\sigma$. We use it for every ``excess over optimum'' number in this paper,
and we never measure progress relative to our own final cost.

\paragraph{Reporting error.} We report the \emph{relative excess}
$C(\sigma)/C^\star-1$ and the \emph{absolute excess per point} $(C(\sigma)-C^\star)/N$ in
cell units squared, where $C^\star$ is a certified optimum. Both are needed.
Inputs that already lie almost on the cells have a tiny $C^\star$, which makes
relative numbers explode. Inputs dominated by outliers (heavy tails) have a huge
$C^\star$, which makes every method look perfect in relative terms.
